% !TEX program = xelatex
\documentclass[a4paper,12pt]{article}

% ===================== 页面与排版 =====================
\usepackage{geometry}
\geometry{left=2.5cm,right=2.5cm,top=2.8cm,bottom=2.8cm}

\usepackage{setspace}
\onehalfspacing
\setlength{\parindent}{2em}
\setlength{\parskip}{0.3em}

% ===================== 数学与物理符号 =====================
\usepackage{amsmath,amssymb}
\usepackage{physics}
\numberwithin{equation}{section}

% ===================== XeLaTeX 中文与字体 =====================
\usepackage{fontspec}
\usepackage{xeCJK}

% 下面字体为 macOS 常见字体；若你系统没有，请替换成你本机存在的字体名
\setmainfont{Times New Roman}
\setmonofont{Courier New}
%\setCJKmainfont{Songti}
%\setCJKsansfont{Heiti}
%\setCJKmonofont{PingFang SC}
%\xeCJKsetup{punctstyle=kaiming}

% ===================== 图形、表格与标题 =====================
\usepackage{graphicx}
\usepackage{float}
\usepackage{subcaption}
\usepackage{booktabs}
\usepackage{caption}
\captionsetup{font=small,labelfont=bf}

% ===================== 代码环境 =====================
\usepackage{listings}
\usepackage{xcolor}

\lstset{
    language=Matlab,
    basicstyle=\ttfamily\small,
    keywordstyle=\color{blue}\bfseries,
    commentstyle=\color{green!40!black}\itshape,
    stringstyle=\color{red!70!black},
    numbers=left,
    numberstyle=\tiny\color{gray},
    frame=single,
    breaklines=true,
    showstringspaces=false,
    captionpos=b
}

% ===================== 页眉页脚 =====================
\usepackage{fancyhdr}
\pagestyle{fancy}
\fancyhf{}
\fancyhead[L]{Numerical Study of Black Hole Spacetimes}
\fancyhead[R]{Tang Yi \& Zhou Jingzi}
\fancyfoot[C]{\thepage}

% ===================== 超链接（尽量最后加载） =====================
\usepackage{hyperref}
\hypersetup{
    colorlinks=true,
    linkcolor=blue,
    citecolor=blue,
    urlcolor=blue
}

% ===================== 标题信息 =====================
\title{\bfseries
广义相对论时空结构的数值模拟与可视化研究\\[0.3em]
\large 从Schwarzschild到RN黑洞}
\author{
学生 \& 同学\\
\small 北京师范大学\ 物理与天文学院\\
\small 学号：XXXXXXXX
}
\date{\small \today}

\begin{document}

\maketitle
\thispagestyle{empty}

\begin{abstract}
本文在广义相对论框架下，面向静态球对称黑洞时空构建了一套可复用的数值模拟与可视化流程，重点比较史瓦西（Schwarzschild）解与赖斯纳--诺德斯特洛姆（Reissner--Nordström, RN）带电黑洞解在因果结构、测地线动力学与光学成像方面的差异。我们实现并串联了若干相互独立的计算模块：彭罗斯图（Penrose diagram）的保角压缩绘制、类时测地线在内视界附近的数值积分、测试粒子轨道进动的有效势分析与轨迹模拟，以及基于逆向光线追踪（backward ray tracing）的黑洞吸积盘成像渲染。数值结果直观展示了RN时空的双视界结构（$r_\pm$）、内视界附近的因果性质变化，以及电荷参数 $Q/M$ 改变对阴影大小与引力透镜特征的影响；特别地，当 $Q/M>1$ 时事件视界消失，对应裸奇点情形，其成像呈现出与常规黑洞显著不同的中心高亮与透镜结构。本文的模拟结果与解析预言一致，为强场区广义相对论的教学与数值实验提供了可视化支撑。
\end{abstract}

\newpage
\tableofcontents
\newpage

% =========================================================
\section{引言与理论框架}

广义相对论（General Relativity, GR）将引力解释为时空几何的体现。其基本动力学由爱因斯坦场方程给出：
\begin{equation}
G_{\mu\nu}\equiv R_{\mu\nu}-\frac{1}{2}Rg_{\mu\nu}=\frac{8\pi G}{c^4}T_{\mu\nu}.
\end{equation}
在四维时空中，度规张量 $g_{\mu\nu}$ 决定线元
\begin{equation}
ds^2=g_{\mu\nu}dx^\mu dx^\nu,
\end{equation}
从而定义了因果结构与测地线（自由落体轨迹）。

\subsection{测地线方程与数值求解}

自由粒子沿测地线运动，满足
\begin{equation}
\frac{d^2x^\mu}{d\tau^2}+\Gamma^\mu_{\alpha\beta}\frac{dx^\alpha}{d\tau}\frac{dx^\beta}{d\tau}=0,
\end{equation}
其中克里斯托费尔符号
\begin{equation}
\Gamma^\mu_{\alpha\beta}
=\frac{1}{2}g^{\mu\sigma}\left(\partial_\beta g_{\sigma\alpha}+\partial_\alpha g_{\sigma\beta}-\partial_\sigma g_{\alpha\beta}\right).
\end{equation}
直接计算 $\Gamma^\mu_{\alpha\beta}$ 容易出错且不利于程序化推广，因此本文主要采用哈密顿表述，将二阶方程转化为一阶正则方程组并进行数值积分。

\subsection{哈密顿表述}

定义正则动量 $p_\mu=g_{\mu\nu}\dot{x}^\nu$（点表示对仿射参数的导数），哈密顿量
\begin{equation}
H=\frac{1}{2}g^{\mu\nu}p_\mu p_\nu.
\end{equation}
对应的正则方程为
\begin{equation}
\frac{dx^\mu}{d\lambda}=\frac{\partial H}{\partial p_\mu},\qquad
\frac{dp_\mu}{d\lambda}=-\frac{\partial H}{\partial x^\mu}.
\end{equation}
在本文的数值实现中，测地线演化、轨道积分与光线追踪均可视为对该一阶系统在不同初值与事件条件下的积分问题。

% =========================================================
\section{静态球对称黑洞度规}

\subsection{史瓦西度规（Schwarzschild metric）}

史瓦西时空描述质量为 $M$ 的不带电静态黑洞。在球坐标 $(t,r,\theta,\phi)$ 下，
\begin{equation}
ds^2=-\left(1-\frac{2M}{r}\right)dt^2+\left(1-\frac{2M}{r}\right)^{-1}dr^2+r^2d\Omega^2,
\end{equation}
其中 $d\Omega^2=d\theta^2+\sin^2\theta\,d\phi^2$。该度规在 $r=2M$ 处出现坐标奇点（事件视界），在 $r=0$ 处为曲率奇点。

常用特征半径包括：
\begin{itemize}
\item 事件视界：$r_s=2M$；
\item 光子球：$r_{\mathrm{ph}}=3M$；
\item 最内稳定圆轨道（ISCO）：$r_{\mathrm{ISCO}}=6M$。
\end{itemize}

\subsection{赖斯纳--诺德斯特洛姆度规（Reissner--Nordström metric）}

当黑洞带电荷 $Q$ 时，电磁场对 $T_{\mu\nu}$ 产生贡献，外部电真空解为
\begin{equation}
ds^2=-f(r)dt^2+\frac{1}{f(r)}dr^2+r^2d\Omega^2,
\end{equation}
其中
\begin{equation}
f(r)=1-\frac{2M}{r}+\frac{Q^2}{r^2}.
\end{equation}
令 $f(r)=0$ 得到两个视界半径
\begin{equation}
r_{\pm}=M\pm\sqrt{M^2-Q^2}.
\end{equation}
其中 $r_+$ 为外视界（event horizon），$r_-$ 为内视界（Cauchy horizon）。

三种典型情形：
\begin{enumerate}
\item $Q<M$：存在内外双视界；
\item $Q=M$：极端黑洞，$r_+=r_-=M$；
\item $Q>M$：无视界，解呈现裸奇点性质。
\end{enumerate}

% =========================================================
\section{彭罗斯图与因果结构}

彭罗斯图通过保角变换将无穷远压缩至有限区域，同时保持光线以 $45^\circ$ 传播，从而清晰呈现时空的因果结构。对静态球对称时空，可先引入乌龟坐标（tortoise coordinate）并构造光锥坐标，再进行保角压缩。

\subsection{坐标变换}

以史瓦西为例，乌龟坐标可写为
\begin{equation}
r^*=r+2M\ln\left|\frac{r}{2M}-1\right|.
\end{equation}
进而定义光锥坐标
\begin{equation}
u=t-r^*,\qquad v=t+r^*,
\end{equation}
并做保角压缩
\begin{equation}
U=\arctan(u),\qquad V=\arctan(v),
\end{equation}
最后取
\begin{equation}
T=\frac{U+V}{2},\qquad R=\frac{V-U}{2}.
\end{equation}

\subsection{史瓦西与RN因果结构对比}

史瓦西黑洞的 $r=0$ 奇点在彭罗斯图上表现为类空（spacelike）边界，即一旦越过事件视界，观察者在有限固有时内不可避免地到达奇点。相比之下，RN时空的奇点在延拓结构中呈现类时（timelike）特征，并存在内视界 $r_-$；这使得其最大解析延拓具有更丰富的区域结构（理论上包含“多宇宙”重复块与类虫洞通道）。需要强调的是，内视界处通常伴随质量膨胀等不稳定性，完整物理可行性需进一步考虑扰动与量子效应；本文主要聚焦经典几何与数值可视化层面。

% =========================================================
\section{测试粒子轨道与近日点进动}

在牛顿引力中，单体引力场下的束缚轨道为闭合椭圆。广义相对论中，时空曲率导致轨道方程包含高阶修正项，从而出现近日点进动。对于RN时空，可通过守恒量与有效势进行分析。

\subsection{有效势与径向方程}

测试粒子径向运动可写为
\begin{equation}
\left(\frac{dr}{d\tau}\right)^2+V_{\mathrm{eff}}(r)=E^2,
\end{equation}
其中（取单位质量）
\begin{equation}
V_{\mathrm{eff}}(r)=\left(1-\frac{2M}{r}+\frac{Q^2}{r^2}\right)\left(1+\frac{L^2}{r^2}\right).
\end{equation}
展开可见包含 $r^{-3}$（相对论修正）与 $r^{-4}$（电荷相关）项，前者是导致近日点进动的重要来源。

\subsection{哈密顿方程的符号推导（示例）}

\begin{lstlisting}[caption={哈密顿方程符号推导示例（RN时空）}]
% 定义符号变量
syms t r th ph pt pr pth pph tau real
syms M Q real

% 度规函数
f_s = 1 - 2*M/r + Q^2/r^2;
g_up = diag([-1/f_s, f_s, 1/r^2, 1/(r^2*sin(th)^2)]);

% 哈密顿量
H_sym = 0.5 * [pt pr pth pph] * g_up * [pt; pr; pth; pph];

% 正则方程
vars = [t; r; th; ph; pt; pr; pth; pph];
eqs = [diff(H_sym, pt); diff(H_sym, pr);
       diff(H_sym, pth); diff(H_sym, pph);
       -diff(H_sym, t); -diff(H_sym, r);
       -diff(H_sym, th); -diff(H_sym, ph)];

odefun = matlabFunction(eqs, 'Vars', {tau, vars});
\end{lstlisting}

\subsection{数值轨道特征}

通过对上述一阶系统进行积分，可以观测到束缚轨道在每圈运动后近日点方向发生偏移。进动量随轨道半径减小而显著增强，与强场区相对论修正增强的直观预期一致。

% =========================================================
\section{逆向光线追踪与黑洞成像}

黑洞阴影与吸积盘成像可通过逆向光线追踪（backward ray tracing）实现：从相机像素出发反向发射光线，沿零测地线积分，判定其是否落入视界、逃向无穷远或与吸积盘相交，并据此为像素赋值。

\subsection{RN度规函数：\texttt{gdown\_RN}}

\begin{equation}
ds^2=-f(r)dt^2+\frac{1}{f(r)}dr^2+r^2(d\theta^2+\sin^2\theta\,d\phi^2),
\quad
f(r)=1-\frac{2M}{r}+\frac{Q^2}{r^2}.
\end{equation}

\begin{lstlisting}[caption={协变度规函数 \texttt{gdown\_RN}}]
function g = gdown_RN(r, th, M, Q)
    % 计算几何因子 f(r)
    f = 1 - 2*M/r + Q^2/r^2;

    % 构造协变度规矩阵 g_{mu nu}
    g = [-f,      0,        0,              0;
          0,    1/f,        0,              0;
          0,      0,      r^2,              0;
          0,      0,        0,  r^2*sin(th)^2];
end
\end{lstlisting}

\subsection{局域正交标架与初始四动量}

为在相机位置将“像素方向”映射为全局四动量，可构造局域正交标架（例如静止观察者的正交基），再将局域四动量投影回坐标基底。该步骤避免直接在曲坐标中手动构造初始条件的不稳定性。

\begin{lstlisting}[caption={局域正交标架 \texttt{localtetrad\_RN}}]
function zamo = localtetrad_RN(metric)
    g11 = metric(1,1);  % g_tt
    g22 = metric(2,2);  % g_rr
    g33 = metric(3,3);  % g_theta theta
    g44 = metric(4,4);  % g_phi phi

    % 正交归一基矢（一个简单的静止标架实现）
    e0 = [1/sqrt(-g11), 0, 0, 0];
    e1 = [0, -1/sqrt(g22), 0, 0];
    e2 = [0, 0, 1/sqrt(g33), 0];
    e3 = [0, 0, 0, -1/sqrt(g44)];

    zamo = [e0; e1; e2; e3];
end
\end{lstlisting}

\begin{lstlisting}[caption={像素到初始四动量映射 \texttt{GetRayMomentum\_RN}}]
function momentumup = GetRayMomentum_RN(metric, fov, npix, i, j)
    % 像素到屏幕坐标
    xs = 2*tan(fov/2)/npix * (i - 0.5*(npix+1));
    ys = 2*tan(fov/2)/npix * (j - 0.5*(npix+1));

    % 屏幕到球面角度
    thx  = atan(sqrt(xs^2 + ys^2));
    psix = atan2(xs, -ys);
    if psix < 0, psix = psix + 2*pi; end

    % 局域空间方向（单位速度方向）
    vel = [cos(thx), sin(thx)*cos(psix), sin(thx)*sin(psix)];

    % 标架
    zamo = localtetrad_RN(metric);

    % 局域四动量（零测地线：可取比例常数 kappa=1）
    kappa = 1;
    p_local = [-kappa, kappa*vel];

    % 投影到坐标基底：p^mu
    momentumup = (p_local * zamo).';
end
\end{lstlisting}

\subsection{测地线方程生成与事件判据}

通过符号工具生成哈密顿方程（避免手算Christoffel符号），并在ODE积分中加入事件检测：落入外视界、逃向远处、或与吸积盘相交。

\begin{lstlisting}[caption={测地线方程生成 \texttt{geodesic\_RN}}]
function odefun = geodesic_RN(M, Q)
    syms t r th ph pt pr pth pph tau real

    f_sym = 1 - 2*M/r + Q^2/r^2;
    g_down_sym = diag([-f_sym, 1/f_sym, r^2, r^2*sin(th)^2]);
    g_up_sym = inv(g_down_sym);

    p = [pt; pr; pth; pph];
    Hsyms = 0.5 * p.' * g_up_sym * p;

    vars = [t; r; th; ph; pt; pr; pth; pph];
    eqs = [diff(Hsyms, pt); diff(Hsyms, pr);
           diff(Hsyms, pth); diff(Hsyms, pph);
           -diff(Hsyms, t); -diff(Hsyms, r);
           -diff(Hsyms, th); -diff(Hsyms, ph)];

    odefun = matlabFunction(eqs, 'Vars', {tau, vars});
end
\end{lstlisting}

\subsection{吸积盘亮度模型（示例）}

可采用薄盘模型的经验标度作为可视化亮度分布，例如
\begin{equation}
T(r)\propto r^{-3/4},\qquad I(r)\propto T^4\propto r^{-3}.
\end{equation}

\begin{lstlisting}[caption={吸积盘亮度函数 \texttt{Diskvalue}}]
function val = Diskvalue(r)
    if r > 6 && r < 20
        val = min(1.2, 10 / r^3);
    else
        val = 0;
    end
end
\end{lstlisting}

\subsection{单光线追踪与着色规则（示例）}

\begin{lstlisting}[caption={单光线追踪 \texttt{TraceSingleRay\_RN}}]
function II = TraceSingleRay_RN(pos, fov, npix, i, j, M, Q, odefun)
    % 初始条件：位置 + 协变动量
    metric = gdown_RN(pos(2), pos(3), M, Q);
    momentumup = GetRayMomentum_RN(metric, fov, npix, i, j);
    momentumdown = metric * momentumup;
    initialconditions = [pos; momentumdown];

    % 外视界
    if M > abs(Q)
        r_horizon = M + sqrt(M^2 - Q^2);
    else
        r_horizon = 0; % 裸奇点：无事件视界
    end

    % 事件检测：入视界 / 逃逸 / 穿过吸积盘面
    function [value, isterminal, direction] = myevent(~, u, r_h)
        value = [u(2) - 1.01*r_h;  % 接近外视界
                 u(2) - 100;       % 远处截断（近似无穷远）
                 u(3) - pi/2];     % 吸积盘平面
        isterminal = [1; 1; 1];
        direction  = [0; 0; 0];
    end

    opts = odeset('Events', @(t,u) myevent(t,u,r_horizon), ...
                  'RelTol', 1e-5, 'AbsTol', 1e-6);

    lambdaf = 500;
    [~, ~, ~, ue, ie] = ode45(odefun, [0 lambdaf], initialconditions, opts);

    % 着色：默认黑
    II = 0;

    if ~isempty(ie)
        last_event = ie(end);

        if last_event == 3
            r_hit = ue(end, 2);
            II = Diskvalue(r_hit);
        elseif last_event == 2
            II = 0.05; % 背景
        else
            II = 0;    % 落入视界
        end
    end
end
\end{lstlisting}

\subsection{不同电荷参数下的成像趋势}

改变 $Q/M$ 会影响 $r_+$、光线偏折强度与阴影尺度。表~\ref{tab:charge} 给出若干代表性参数下的定性对比（数值可据实际仿真输出更新）。

\begin{table}[H]
\centering
\caption{不同电荷下的黑洞特性（示例表，建议用你的仿真输出替换）}
\label{tab:charge}
\begin{tabular}{cccc}
\toprule
$Q/M$ & 外视界 $r_+/M$ & 阴影角半径（示意） & 物理状态 \\
\midrule
0.0  & 2.00 & $3.46^\circ$ & Schwarzschild \\
0.5  & 1.87 & $3.34^\circ$ & RN（双视界） \\
0.75 & 1.66 & $3.15^\circ$ & RN（双视界） \\
1.0  & 1.00 & $2.32^\circ$ & 极端RN \\
1.15 & --   & --           & 裸奇点（无视界） \\
\bottomrule
\end{tabular}
\end{table}

% =========================================================
\section{结论与展望}

本文构建并实现了一套用于静态球对称黑洞时空的数值模拟与可视化流程，涵盖因果结构（彭罗斯图）、测地线动力学（穿越视界与内视界附近演化）、轨道进动（有效势与轨迹积分）以及成像问题（逆向光线追踪与吸积盘渲染）。数值结果与史瓦西与RN解的解析性质一致，并直观呈现了电荷参数对视界结构与光学观测量（阴影尺度、透镜结构）的系统性影响。对于 $Q/M>1$ 的无视界情形，成像结果出现与常规黑洞显著不同的结构特征，为理解经典时空解在不同参数区间的几何与观测差异提供了直接的数值证据。

后续工作可沿两个方向扩展：其一，将静态带电推广至旋转带电的 Kerr--Newman 时空；其二，引入更真实的辐射传输与盘模型（速度场、多普勒增亮、光学厚度等），以获得更接近天体物理观测的合成图像。

% =========================================================
\section*{致谢}
感谢导师的指导与讨论。本文的数值实现参考了广义相对论与黑洞物理的经典教材与文献。

% =========================================================
\begin{thebibliography}{99}

\bibitem{mtw}
C.~W.~Misner, K.~S.~Thorne, and J.~A.~Wheeler,
\emph{Gravitation},
W. H. Freeman (1973).

\bibitem{wald}
R.~M.~Wald,
\emph{General Relativity},
University of Chicago Press (1984).

\bibitem{chandrasekhar}
S.~Chandrasekhar,
\emph{The Mathematical Theory of Black Holes},
Oxford University Press (1983).

\bibitem{shapiro}
S.~L.~Shapiro and S.~A.~Teukolsky,
\emph{Black Holes, White Dwarfs, and Neutron Stars},
Wiley (1983).

\bibitem{penrose}
R.~Penrose,
``Gravitational Collapse and Space-Time Singularities'',
\emph{Phys. Rev. Lett.} \textbf{14}, 57 (1965).

\bibitem{eht}
Event Horizon Telescope Collaboration,
``First M87 Event Horizon Telescope Results'',
\emph{Astrophys. J. Lett.} \textbf{875}, L1 (2019).

\end{thebibliography}

\end{document}
